% OpenStax Calculus Volume 2 -- Section 6.1 Exercises: Power Series and Functions
% Exercises reproduced from OpenStax, licensed under CC BY 4.0.
% Source: https://openstax.org/books/calculus-volume-2/pages/6-1-power-series-and-functions
\documentclass{exercises}
\title{OpenStax Calculus Volume 2}
\subtitle{Section 6.1 Exercises: Power Series and Functions}
\sourceurl{https://openstax.org/books/calculus-volume-2/pages/6-1-power-series-and-functions}
\author{}
\date{}

\begin{document}
\maketitle


In the following exercises, state whether each statement is true, or give an example to show that it is false.

\begin{enumerate}[start=1]
\item If $\sum_{n=1}^{\infty}a_{n}x^{n}$ converges, then $a_{n}x^{n}\to0$ as $n\to \infty$.
\item $\sum_{n=1}^{\infty}a_{n}x^{n}$ converges at $x=0$ for any real numbers $a_{n}$.
\item Given any sequence $a_{n}$, there is always some $R>0$, possibly very small, such that $\sum_{n=1}^{\infty}a_{n}x^{n}$ converges on $(-R,R)$.
\item If $\sum_{n=1}^{\infty}a_{n}x^{n}$ has radius of convergence $R>0$ and if $|b_{n}|\le|a_{n}|$ for all $n$, then the radius of convergence of $\sum_{n=1}^{\infty}b_{n}x^{n}$ is greater than or equal to $R$.
\item Suppose that $\sum_{n=0}^{\infty}a_{n}{(x-3)}^{n}$ converges at $x=6$. At which of the following points might the series diverge? Use the fact that if $\sum a_{n}{(x-c)}^{n}$ converges at $x$, then it converges at any point closer to $c$ than $x$.
\begin{enumerate}[label=(\alph*)]
  \item $x=1$
  \item $x=2$
  \item $x=3$
  \item $x=0$
  \item $x=5.99$
  \item $x=0.000001$
\end{enumerate}
\item Suppose that $\sum_{n=0}^{\infty}a_{n}{(x+1)}^{n}$ converges at $x=-2$. At which of the following points must the series also converge? Use the fact that if $\sum a_{n}{(x-c)}^{n}$ converges at $x$, then it converges at any point closer to $c$ than $x$.
\begin{enumerate}[label=(\alph*)]
  \item $x=2$
  \item $x=-1$
  \item $x=-3$
  \item $x=0$
  \item $x=0.99$
  \item $x=0.000001$
\end{enumerate}
\end{enumerate}

In the following exercises, suppose that $|\frac{a_{n+1}}{a_{n}}|\to1$ as $n\to \infty$. Find the radius of convergence for each series.

\begin{enumerate}[resume]
\item $\sum_{n=0}^{\infty}a_{n}2^{n}x^{n}$
\item $\sum_{n=0}^{\infty}\frac{a_{n}x^{n}}{2^{n}}$
\item $\sum_{n=0}^{\infty}\frac{a_{n}\pi^{n}x^{n}}{e^{n}}$
\item $\sum_{n=0}^{\infty}\frac{a_{n}{(-1)}^{n}x^{n}}{10^{n}}$
\item $\sum_{n=0}^{\infty}a_{n}{(-1)}^{n}x^{2n}$
\item $\sum_{n=0}^{\infty}a_{n}{(-4)}^{n}x^{2n}$
\end{enumerate}

In the following exercises, find the radius of convergence $R$ and interval of convergence for $\sum a_{n}x^{n}$ with the given coefficients $a_{n}$.

\begin{enumerate}[resume]
\item $\sum_{n=1}^{\infty}\frac{{(2x)}^{n}}{n}$
\item $\sum_{n=1}^{\infty}{(-1)}^{n}\frac{x^{n}}{\sqrt{n}}$
\item $\sum_{n=1}^{\infty}\frac{nx^{n}}{2^{n}}$
\item $\sum_{n=1}^{\infty}\frac{nx^{n}}{e^{n}}$
\item $\sum_{n=1}^{\infty}\frac{n^{2}x^{n}}{2^{n}}$
\item $\sum_{k=1}^{\infty}\frac{k^{e}x^{k}}{e^{k}}$
\item $\sum_{k=1}^{\infty}\frac{\pi^{k}x^{k}}{k^{\pi}}$
\item $\sum_{n=1}^{\infty}\frac{x^{n}}{n!}$
\item $\sum_{n=1}^{\infty}\frac{10^{n}x^{n}}{n!}$
\item $\sum_{n=1}^{\infty}{(-1)}^{n}\frac{x^{n}}{\ln (2n)}$
\end{enumerate}

In the following exercises, find the radius of convergence of each series.

\begin{enumerate}[resume]
\item $\sum_{k=1}^{\infty}\frac{{(k!)}^{2}x^{k}}{(2k)!}$
\item $\sum_{n=1}^{\infty}\frac{(2n)!x^{n}}{n^{2n}}$
\item $\sum_{k=1}^{\infty}\frac{k!}{1\cdot3\cdot5\cdots(2k-1)}x^{k}$
\item $\sum_{k=1}^{\infty}\frac{2\cdot4\cdot6\cdots2k}{(2k)!}x^{k}$
\item $\sum_{n=1}^{\infty}\frac{x^{n}}{\binom{2n}{n}}$ where $\binom{n}{k}=\frac{n!}{k!(n-k)!}$
\item $\sum_{n=1}^{\infty}\sin^{2}(n)x^{n}$
\end{enumerate}

In the following exercises, use the ratio test to determine the radius of convergence of each series.

\begin{enumerate}[resume]
\item $\sum_{n=1}^{\infty}\frac{{(n!)}^{3}}{(3n)!}x^{n}$
\item $\sum_{n=1}^{\infty}\frac{2^{3n}{(n!)}^{3}}{(3n)!}x^{n}$
\item $\sum_{n=1}^{\infty}\frac{n!}{n^{n}}x^{n}$
\item $\sum_{n=1}^{\infty}\frac{(2n)!}{n^{2n}}x^{n}$
\end{enumerate}

In the following exercises, given that $\frac{1}{1-x}=\sum_{n=0}^{\infty}x^{n}$ with convergence in $(-1,1)$, find the power series for each function with the given center $a$, and identify its interval of convergence.

\begin{enumerate}[resume]
\item $f(x)=\frac{1}{x}$; $a=1$. (Hint: $\frac{1}{x}=\frac{1}{1-(1-x)}$.)
\item $f(x)=\frac{1}{1-x^{2}}$, $a=0$
\item $f(x)=\frac{x}{1-x^{2}}$, $a=0$
\item $f(x)=\frac{1}{1+x^{2}}$, $a=0$
\item $f(x)=\frac{x^{2}}{1+x^{2}}$, $a=0$
\item $f(x)=\frac{1}{2-x}$, $a=1$
\item $f(x)=\frac{1}{1-2x}$; $a=0$.
\item $f(x)=\frac{1}{1-4x^{2}}$, $a=0$
\item $f(x)=\frac{x^{2}}{1-4x^{2}}$, $a=0$
\item $f(x)=\frac{x^{2}}{5-4x+x^{2}}$, $a=2$
\end{enumerate}

Use the next exercise to find the radius of convergence of the given series in the subsequent exercises.

\begin{enumerate}[resume]
\item Explain why, if ${|a_{n}|}^{1/n}\to r>0$, then ${|a_{n}x^{n}|}^{1/n}\to|x|r<1$ whenever $|x|<\frac{1}{r}$ and, therefore, the radius of convergence of $\sum_{n=1}^{\infty}a_{n}x^{n}$ is $R=\frac{1}{r}$.
\item $\sum_{n=1}^{\infty}\frac{x^{n}}{n^{n}}$
\item $\sum_{k=1}^{\infty}{(\frac{k-1}{2k+3})}^{k}x^{k}$
\item $\sum_{k=1}^{\infty}{(\frac{2k^{2}-1}{k^{2}+3})}^{k}x^{k}$
\item $\sum_{n=1}^{\infty}a_{n}x^{n}$, where $a_{n}={(n^{1/n}-1)}^{n}$
\item Suppose that $p(x)=\sum_{n=0}^{\infty}a_{n}x^{n}$ such that $a_{n}=0$ if $n$ is even. Explain why $p(x)=-p(-x)$.
\item Suppose that $p(x)=\sum_{n=0}^{\infty}a_{n}x^{n}$ such that $a_{n}=0$ if $n$ is odd. Explain why $p(x)=p(-x)$.
\item Suppose that $p(x)=\sum_{n=0}^{\infty}a_{n}x^{n}$ converges on $(-1,1]$. Find the interval of convergence of $p(Ax)$.
\item Suppose that $p(x)=\sum_{n=0}^{\infty}a_{n}x^{n}$ converges on $(-1,1]$. Find the interval of convergence of $p(2x-1)$.
\end{enumerate}

In the following exercises, suppose that $p(x)=\sum_{n=0}^{\infty}a_{n}x^{n}$ satisfies $\lim_{n\to \infty}\frac{a_{n+1}}{a_{n}}=1$ where $a_{n}\ge0$ for each $n$. State whether each series converges on the full interval $(-1,1)$, or if there is not enough information to draw a conclusion. Use the comparison test when appropriate.

\begin{enumerate}[resume]
\item $\sum_{n=0}^{\infty}a_{n}x^{2n}$
\item $\sum_{n=0}^{\infty}a_{2n}x^{2n}$
\item $\sum_{n=0}^{\infty}a_{2n}x^{n}$ (Hint: $x=\pm\sqrt{x^{2}}$.)
\item $\sum_{n=0}^{\infty}a_{n^{2}}x^{n^{2}}$ (Hint: Let $b_{k}=a_{k}$ if $k=n^{2}$ for some $n$; otherwise, $b_{k}=0$.)
\item Suppose that $p(x)$ is a polynomial of degree $N$. Find the radius and interval of convergence of $\sum_{n=1}^{\infty}p(n)x^{n}$.
\item \techrequired Plot the graphs of $\frac{1}{1-x}$ and of the partial sums $S_{N}=\sum_{n=0}^{N}x^{n}$ for $N=10,20,30$ on the interval $[-0.99,0.99]$. Comment on the approximation of $\frac{1}{1-x}$ by $S_{N}$ near $x=-1$ and near $x=1$ as $N$ increases.
\item \techrequired Plot the graphs of $-\ln (1-x)$ and of the partial sums $S_{N}=\sum_{n=1}^{N}\frac{x^{n}}{n}$ for $N=10,50,100$ on the interval $[-0.99,0.99]$. Comment on the behavior of the sums near $x=-1$ and near $x=1$ as $N$ increases.
\item \techrequired Plot the graphs of the partial sums $S_{N}=\sum_{n=1}^{N}\frac{x^{n}}{n^{2}}$ for $N=10,50,100$ on the interval $[-0.99,0.99]$. Comment on the behavior of the sums near $x=-1$ and near $x=1$ as $N$ increases.
\item \techrequired Plot the graphs of the partial sums $S_{N}=\sum_{n=1}^{N}\sin(n)x^{n}$ for $N=10,50,100$ on the interval $[-0.99,0.99]$. Comment on the behavior of the sums near $x=-1$ and near $x=1$ as $N$ increases.
\item \techrequired Plot the graphs of the partial sums $S_{N}=\sum_{n=0}^{N}{(-1)}^{n}\frac{x^{2n+1}}{(2n+1)!}$ for $N=3,5,10$ on the interval $[-2\pi,2\pi]$. Comment on how these plots approximate $\sin x$ as $N$ increases.
\item \techrequired Plot the graphs of the partial sums $S_{N}=\sum_{n=0}^{N}{(-1)}^{n}\frac{x^{2n}}{(2n)!}$ for $N=3,5,10$ on the interval $[-2\pi,2\pi]$. Comment on how these plots approximate $\cos x$ as $N$ increases.
\end{enumerate}

\end{document}
